\documentclass[10pt,letterpaper]{article}

\usepackage[margin=0.48in]{geometry}
\usepackage[T1]{fontenc}
\usepackage{lmodern}
\usepackage{xcolor}
\usepackage{enumitem}
\usepackage{tabularx}
\usepackage[hidelinks]{hyperref}
\usepackage{titlesec}
\usepackage{microtype}

\definecolor{accent}{HTML}{176B70}
\setlength{\parindent}{0pt}
\setlength{\parskip}{0pt}
\setlist[itemize]{leftmargin=1.15em,itemsep=0.7pt,topsep=1pt,parsep=0pt}
\titleformat{\section}{\large\bfseries\color{accent}}{}{0pt}{}[\vspace{-2pt}\titlerule]
\titlespacing*{\section}{0pt}{4pt}{2pt}
\pagestyle{empty}

\newcommand{\role}[4]{%
  \textbf{#1}\hfill #2\\
  \textit{#3}\hfill \textit{#4}\vspace{-3pt}
}

\begin{document}
\fontsize{9}{10.4}\selectfont

\begin{center}
  {\Large\bfseries Piter Z. Garcia Bautista}\\[1pt]
  {\large Healthcare Data, Business Analysis, and Applied AI}\\[2pt]
  Rochester, NY \enspace|\enspace Remote now; open to relocation after December 2026\\
  \href{mailto:pzg8794@rit.edu}{pzg8794@rit.edu}
  \enspace|\enspace +1 (646) 288-3300
  \enspace|\enspace \href{https://linkedin.com/in/garciapiterz}{linkedin.com/in/garciapiterz}\\
  \href{https://github.com/pzg8794}{github.com/pzg8794}
  \enspace|\enspace \href{https://garciapiterz.com}{garciapiterz.com}
\end{center}

\section*{Summary}
Graduate data scientist and healthcare AI practitioner with experience
building Python analytics workflows, validating production data, documenting
requirements, and translating technical results for varied stakeholders.
Current work spans applied machine learning, healthcare analytics,
bioinformatics, reproducible evaluation, and inclusive technical
communication. Seeking to support clinical AI implementation through careful
data analysis, process documentation, testing, and cross-functional
collaboration.

\section*{Relevant Skills}
\begin{tabularx}{\textwidth}{@{}>{\bfseries}l X@{}}
Data analysis & Python, SQL, R, pandas, NumPy, SciPy, scikit-learn, structured data review, statistical analysis, data-quality validation\\
Business analysis & Requirements clarification, workflow documentation, stakeholder communication, acceptance-oriented testing, defect and result reporting\\
Healthcare and AI & Healthcare machine learning, hospital-readmission analytics, bioinformatics, clinical-data research, responsible and fairness-aware AI\\
Delivery & AWS, GCP, Docker, Git/GitHub, Linux/Unix, reproducible pipelines, structured logging, technical reports and presentations\\
\end{tabularx}

\section*{Relevant Experience}
\role{Data Solutions Engineer}{Sep 2022--Aug 2024}
{VEDADATA Inc.}{New York, NY}
\begin{itemize}
  \item Designed maintainable Python and pandas workflows for data ingestion,
    cleaning, preprocessing, validation, and analytics-ready delivery.
  \item Built quality checks and statistical diagnostics to identify missing,
    inconsistent, and unusual information in production data workflows.
  \item Collaborated across technical and nontechnical teams to clarify
    requirements, review outputs, document findings, and improve system
    reliability and usability.
\end{itemize}

\role{AI Data Solutions Engineer}{Jan 2021--Sep 2022}
{VIOME Inc.}{New York, NY}
\begin{itemize}
  \item Developed healthcare-oriented machine-learning workflows for
    preprocessing, feature engineering, supervised-model experimentation, and
    evaluation.
  \item Supported reproducible analytics and deployment-oriented data
    preparation in AWS environments for predictive health applications.
\end{itemize}

\role{Graduate Assistant, Quantum and AI Research}{Fall 2025--Present}
{Rochester Institute of Technology, Software Engineering}{Rochester, NY}
\begin{itemize}
  \item Build Python evaluation workflows for adaptive decision systems,
    including automated validation, structured logging, shared datasets, and
    reproducible experiment tracking.
  \item Translate complex requirements and large-scale results into benchmark
    analyses, technical reports, presentations, and manuscript-ready artifacts.
\end{itemize}

\role{Computer Science Teacher Candidate}{2025--Present}
{University of Rochester / Pine Brook Elementary School}{Rochester, NY}
\begin{itemize}
  \item Co-design accessible technical learning and communicate complex ideas
    through multiple formats for learners and collaborators with varied needs.
\end{itemize}

\section*{Selected Healthcare and Data Projects}
\textbf{Hospital Readmission Analytics.}
Compared ensemble and regularized regression methods for clinical risk
stratification, with structured preprocessing, evaluation, and interpretation.

\textbf{Fairness-Aware Bioinformatics and Healthcare Analytics.}
Connected RNA-focused computational workflows, predictive modeling, and
fairness assessment to evaluate technical performance and equitable outcomes.

\textbf{Reproducible Data-System Evaluation.}
Built validation and reporting workflows that convert heterogeneous experiment
outputs into comparison-ready datasets, diagnostics, and stakeholder-facing
artifacts.

\section*{Education}
\textbf{M.S., Data Science}, Rochester Institute of Technology
\hfill Expected December 2026\\
GPA: 3.9+; applied AI, bioinformatics, neural networks, and reproducible data systems

\textbf{M.S., Teaching Computer Science K--12}, University of Rochester
\hfill Expected December 2026\\
GPA: 3.9+; NSF Robert Noyce Scholar; Advanced Certificate in Leadership in
Disability and Inclusive Practices

\textbf{M.S., Computer Science}, Rochester Institute of Technology
\hfill 2015\\
Artificial intelligence, big data analytics, and computer graphics

\textbf{B.S., Computer Engineering Technology}, Farmingdale State College
\hfill 2012

\end{document}
